\section{综合检验、适用范围与防灾策略}
\subsection{结果解释与适用范围}
各模型回答的对象不同，主要检验及其不能推出的结论见表\ref{tab:validation}。统计判别、数值闭合和条件成员分位不能自动跨越到真实灾害发生概率或现场预测技巧。

\begin{table}[htbp]\centering\small
\caption{检验内容与适用范围}\label{tab:validation}
\begin{tabularx}{\textwidth}{Y Y Y}\toprule
检验&支持内容&不能推出的结论\\\midrule
五层静力订正与同压对照&44018个共同小时中9016个稳定度分类改变&参考高度就是真实安装高度\\
31项统一候选、相关簇与留年检查&固定候选与同季匹配下的事件前环境异常排序&真实灾害因果、发生概率或普适前兆\\
完整床面受力与78/6状态分支&给定床层参数下的水流门槛和重力失稳分类&所有沟道的现场稳定性或安全认证\\
三组IFS起报的3000配对成员&固定08:00发布时可形成条件源区与即时释放到达分位&真实泥石流形成区间或频率校准\\
10:30单点与融水提示核验&辅助即时释放区间包含1/1、提示时差2.5小时&真实到达覆盖、真实灾害预警或长期覆盖率\\
确定性与结果一致性检查&单位、状态、支持与来源一致&自然过程被物理验证\\
D8质量收支、CFL与终止状态&固定地形和给定参数下的离散守恒与右删失状态&自然最大运程或实测危险区\\
一维自然停止与D8到达&两种结构在同一事件下的显著分歧&可从两引擎拼接最优时间或距离\\
\bottomrule
\end{tabularx}\end{table}

这些负结果分别限定了前兆、水力门槛、真实预警和自然运程的外推：问题一前五在留一年检验中仅重合3--4项，问题二峰值触发比0.070032未识别冰崩链，问题三真实释放与到达不可估计，问题四D8全部右删失而一维自然停止。辅助区间包含10:30，不代表真实覆盖。

\subsection{不确定性的来源}
观测与语义不确定度包括站点代表性、质量码释义、附件4冰面/冰下语义、附件7采集时刻、口岸断面吸附和历史长度定义。参数不确定度包括曲线数、RGI6编目面积权重、释放比例、物源放大、床厚与摩擦角、Manning糙率、屈服门槛、D8配置和有限库存幅度。结构不确定度包括静力积分假设、集总产流、融水代理、AR(1)残差、相关簇阈值、单流向路由和右删失分析域。

IFS成员内部同时扰动回归系数、连续小时残差和传播配置，但三者不是由多事件联合估计，抽样概率只属于给定模型集合\cite{beven1992}。D8十六个配置等权也不表示现实配置概率；改变集合会使条件分位移动。质量残差$2.075820\times10^{-15}$和最大Courant数0.188281说明数值实现满足给定约束，不说明地形时效、源区体积、屈服门槛或流变正确。

\subsection{分层预警与监测建议}
第一层是固定时次的环境值守。可在每日固定时次计算31项关联指标、近地层稳定度和水源比，但只有单事件关联的头部特征应作为背景解释，不能直接升级为灾害概率。

第二层是源区监测与前向预报。优先补充冰体位移、裂缝、地震或声学释放信号以及可核查的泥石流形成时刻；将真实$t_{\mathrm{trig}}$事件目录用于重新训练、区间校准和独立覆盖，而不是用$t_m$或口岸冲击替代。发布端应固定数据截止规则，记录模式起报、实际可用延迟和航线；本次六小时分发上界与10:30锚点只构成2.5小时融水提示时差，不足以保证真实灾害预警获得正提前量。

第三层是传播与现场处置。D8到达时间只适用于同一体积、断面、地形和参数设定；若物质前锋进入分析域末端，应继续扩展域或标记右删失，不把26045.494670米当作自然边界。沿程物源节点需有库存和可动性依据才能进入主情景；堰塞湖漫溢作为背景风险，溃决流量需另建坝体、库容和出流模型。

\subsection{模型评价与改进顺序}
本文的改进在于结果可追溯，而非追求表面精度：前五不排除API与净辐射，床面门槛含沿坡自重并区分重力失稳，问题三固定发布时次，问题四区分自然停止与域删失；代价是保留多处不利结论，并维护结果与数据来源的对应关系。

后续改进按观测收益排序：先获取真实形成时刻、源区释放和现场流深流速；其次补充床厚、材料强度、孔压及沿程库存；再以独立多事件目录训练和检验时间模型；最后才考虑二维横向扩散或复杂流变。扩大候选池、挑选有利端点或放宽区间都不能替代这些观测。

\section{结论}
本文形成共享候选、几何、门槛、发布时刻和传播参数的条件分析框架。
\begin{enumerate}
\item 五层静力订正使9016/44018个可算小时的稳定度分类改变，占20.48\%；订正后$Ri_b$中位$-0.0294$，夜间逆温频率60.07\%。31项统一候选去冗余后的前五为API、净辐射、24小时升温、0.1米土壤含水率和潜热通量，$A$为3.829043至1.133226；该结果限定为本次事件前环境关联。
\item 床面受力纳入床体沿坡自重后，84条历史沟道分为78条水流控制和6条重力失稳。正式中心线$q_c=67.588$ m$^2$/s、$Q_c=2.893667\times10^3$ m$^3$/s；事件窗水源峰值202.650 m$^3$/s，最大$T=0.070032$，未出现$T\ge1$。335个源区候选像元不构成冰体失稳识别，严格共同域空间验证不可估计。
\item 固定08:00发布的三组IFS起报生成3000个有限配对成员。源区融水代理区间为09:03:22.847--10:57:51.047，即时释放条件到达区间为09:35:26.068--11:31:59.410，$p_{\mathrm{actionable}}=1$；真实释放、形成和实际到达覆盖不可估计，真实灾害预警尚不能给出。10:30落入辅助即时释放区间（1/1），但辅助到达区间提前量为$-1.033169$至$+0.909426$小时，信息时差2.5小时不构成命中提前量；同链沿沟距离22.924千米，16配置首达17.570--44.728分钟。
\item 问题四的混合入沟体积为$2.635835\times10^6$ m$^3$。D8基线32.766603分钟到达口岸代理，质量相对残差$2.075820\times10^{-15}$、最大Courant数0.188281；16个配置均在26045.494670米分析域末端右删失，不足以外推自然运程。一维基线9.336769千米自然停止；144个无节点配置中48个正速通过、96个首次停止。
\item ERA5与MERRA2同日降水为7.50和1.22毫米；IFS起报时刻与产品版本已分类保存。真实形成时刻、现场传播观测和沿程库存仍缺失，故本文给出条件预报及其适用范围，不提供真实灾害95\%覆盖或工程设防保证。
\end{enumerate}
